\scanpage{124}
i.e.
\[
(\lambda-1)q=\pi(q(c)\pi+\psi(c))\qquad
\left(\begin{aligned}\psi(c)&\in\Gamma(S,\check T)\\\psi(c)(x)&=\varphi(c,x)\end{aligned}\right)
\]
Comme $q$ est \deleted{non dégénérée}, lisse en tte fibre, ceci montre, si $\operatorname{rg}T\geq3$, i.e. $\operatorname{rg}E\geq2$, que $\lambda-1$ est nul en tte fibre (i.e que $\lambda-1$ doit être nilpotent), \deleted{\unsure{on}} et \deleted{que} comme $\pi$ est $\ne0$ en tte fibre, qu'on doit avoir $\psi(c)+q(c)\pi=0$ en tte fibre, \ldots\ Regardons les solutions correspondant à $\lambda=1$, i.e.
\[
\unsure{(|C}\qquad q\circ\sigma=q
\]
qui équivaut à
\[
q(c)\pi+\psi(c)=0\tag*{(19)}
\]
\deleted{On peut} On peut prendre p.ex. $c=-a$ (dans le cas \unsure{$T$|$E$} de rg pair, qui nous intéresse surtout)
\[
\sigma(x)=x\unsure{-|+}\pi(x)a\qquad(\operatorname{rg}\unsure{T|E}\ \text{pair})\tag*{(20)}
\]
% margin note written sideways
\margin{\uline{Prop} Il existe une unique réflexion de $P$ admettant $H$ comme hyperplan des pts fixes et \unsure{invariant|laissant} $Q$. Sa restriction à $\unsure{T|E}$ est donnée par (20) $\sigma x=x-\pi(x)a$ i.e. $\sigma=\mathrm{id}-\pi\otimes a$}
qui est de plus une réflexion i.e.
\[
\sigma^2=\mathrm{Id}_T\tag*{(21)}
\]
(car \deleted{\unsure{$\pi=$|$\sigma^2=$}} $\pi(a)=2\ldots$) En fait, c'est la seule solution \inserted{partt $\ne\mathrm{id}$}, \deleted{\unsure{on|dans}} car de (19) on tire $\psi(c)=\lambda\pi$ ($\lambda=-q(c)$) \deleted{\unsure{d'où}} i.e. $\psi(c)=\lambda\psi(a)=\psi(\lambda a)$ i.e. (comme $\psi$ est iso) $c=\underline{\lambda a}$, et il faut $\lambda=-q(c)$ i.e. $\lambda=-\lambda^2$ i.e. $\lambda(\lambda+1)=0$, ce qui (si on veut $\lambda$ partout non nul) donne $\lambda=-1$\ldots.

On dit que $\sigma$ est la \uline{symétrie canonique} de la quadrique affine régulière.
% archive sheet number: 123
